We formulate the Rivero--de Vries electroweak observation as a two-branch secular problem built from a Casimir variable \(J=s(s+1)\).  The positive root of \(x^2+Jx-J=0\), evaluated at \(J=3/4\) and \(J=2\), gives \(1-x_+(3/4)/x_+(2)=0.22310132\ldots\), close to the weak angle defined by the W/Z pole-mass ratio.  Under the vector-spectrum normalization, the same secular equation has a negative branch that supplies a candidate scalar/order-parameter datum; this scalar reading is conjectural until a gauge-invariant map is derived.  The paper treats the two-branch pattern as a conceptual clue about where a string/Kaluza--Klein spectral datum might attach to the Standard Model.  The proposed attachment point is the low-energy pole spectrum.  The comparison point is the grand-unified value \(\sin^2\theta_W=3/8\), a high-scale group-theoretic datum subsequently related to experiment by renormalization-group flow.  The paper returns to the early dual-model and Kaluza--Klein ambition that string theory should organize observed particle structure, and keeps the modern distinctions among pole masses, running parameters, compactification spectra, and symmetry-breaking data.  We present the algebraic construction, the electroweak projective interpretation in which the Higgs vacuum scale is radial, and the analytical tasks required for a dynamical derivation.  The central open calculation is the derivation of the ordered assignment \((J_W,J_Z)=(3/4,2)\) from the gauge-Higgs/order-parameter sector or from a compactification boundary problem.
